\chapter{Разработка навигационного алгоритма счисления координат и ориентации}
\label{chap:navigation}

% -----------------------------------------------------------------------------

Для определения пространственного перемещения путеизмерительной тележки по показаниям микромеханических датчиков недостаточно располагать только исходными измерениями акселерометра и гироскопа. Необходимо реализовать алгоритм, который преобразует эти измерения в траекторию движения, учитывая физические особенности инерциальной навигации: накопление ошибок при численном интегрировании, влияние гравитации, проекцию угловой скорости вращения Земли на оси чувствительности датчиков и инструментальные погрешности самих сенсоров.

Целью данной главы является разработка навигационного алгоритма счисления координат и ориентации путеизмерительной тележки, адаптированного к условиям её эксплуатации. Основные задачи, решаемые в главе:

\begin{enumerate}
	\item обоснование выбора системы координат, в которой описывается движение тележки;
	\item построение математического аппарата для описания ориентации тела и счисления координат на основе кватернионов;
	\item реализация процедуры начальной выставки (инициализации) с использованием статического буфера перед началом движения, включающей оценку начальной ориентации тела относительно навигационной системы и вычисление смещений нуля датчиков;
	\item построение траектории движения путём численного интегрирования угловой скорости и ускорения с введением коррекции накопленных ошибок на основе априорной информации.
\end{enumerate}

Особенностью предлагаемого подхода является учёт вращения Земли при начальной выставке и в процессе счисления пути, что позволяет минимизировать ошибки, связанные с проекцией угловой скорости Земли на оси чувствительности гироскопа. Для представления ориентации используются кватернионы, обладающие преимуществами перед углами Эйлера и матрицами направляющих косинусов с точки зрения вычислительной устойчивости и отсутствия сингулярностей. Коррекция конечных условий осуществляется по линейной модели, предполагающей доминирование постоянной составляющей ошибки ускорения.

% -----------------------------------------------------------------------------

\section{Выбор системы координат}
\label{sec:coord_systems}

Для математического описания движения путеизмерительной тележки необходимо ввести две правые ортогональные системы координат: связанную с тележкой систему координат (body frame, далее $B$) и навигационную систему координат (navigation frame, далее $N$), относительно которой определяется траектория.

Связанная система координат $B$ жёстко закреплена в корпусе путеизмерительной тележки. Её оси совпадают с осями чувствительности МЕМС-датчиков, расположенных на печатной плате микроконтроллера.

В качестве навигационной выбрана система \textbf{ENU} (East–North–Up). Её оси ориентированы следующим образом:
\begin{enumerate}
	\item $OX_{N}$ - на восток (East);
	\item $OY_{N}$ - на север (North);
	\item $OZ_{N}$ - по вертикали (Up);
\end{enumerate}

Выбор ENU обусловлен удобством интерпретации результатов измерений (профиль пути в вертикальной плоскости, план в горизонтальной) и широким распространением в задачах инерциальной навигации наземных объектов. Кроме того, в этой системе просто записываются векторы гравитации и угловой скорости вращения Земли, что важно для процедур начальной выставки.

Задача навигационного алгоритма заключается в определении ориентации системы $B$ относительно $N$ в каждый момент времени и последующем вычислении траектории.

% -----------------------------------------------------------------------------

\section{Описание математического аппарата для расчёта траектории движения}
\label{sec:math_apparatus}

Для описания взаимной ориентации связанной системы $B$ и навигационной системы $N$ в пространстве могут быть использованы три основных математических аппарата: углы Эйлера, матрица направляющих косинусов и кватернионы. Кратко рассмотрим их достоинства и недостатки применительно к задаче инерциального счисления пути.

\subsection{Углы Эйлера}
Углы Эйлера (тангаж $\theta$, крен $\gamma$, курс $\psi$) интуитивно понятны и имеют ясную физическую интерпретацию. Однако их использование в навигационных алгоритмах сопряжено с двумя фундаментальными проблемами:

\begin{enumerate}[label=\arabic*]
	\item Наличие сингулярностей: при определённых значениях углов (например, $\theta = \pm \pi / 2$) происходит вырождение системы кинематических уравнений – так называемый «гимбал-лок» (gimbal lock);
	\item Вычислительная сложность: уравнения содержат тригонометрические функции, что увеличивает вычислительную нагрузку.
\end{enumerate}


\subsection{Матрица направляющих косинусов}
Матрица направляющих косинусов размера $3 \times 3$ не имеет сингулярностей, а её элементы непосредственно используются для поворота векторов. Однако матрица содержит 9 элементов, хотя реальную ориентацию определяют лишь 3 независимых параметра. Это приводит к избыточности и необходимости поддержания ортогональности (ортогонализации) на каждом шаге интегрирования.


\subsection{Кватернионы}
Кватернионы представляют собой систему гиперкомплексных чисел с одной действительной и тремя мнимыми частями, которая образует векторное пространство над полем вещественных чисел.\\

Кватернионная форма описания ориентации обладает рядом преимуществ:
\begin{enumerate}
	\item Отсутствие сингулярностей. Кватернионы свободны от проблемы «гимбал-лока» при любых углах поворота, что гарантирует работоспособность алгоритма во всех возможных ориентациях тележки;
	\item Компактность и вычислительная эффективность. Для хранения и обработки требуется всего 4 числа вместо 9 у матрицы. Операции умножения кватернионов и поворота вектора требуют меньшего количества арифметических операций, что повышает вычислительную эффективность алгоритма;
	\item Простота нормализации. В процессе численного интегрирования ошибки округления могут приводить к отклонению нормы кватерниона от единицы. Восстановление условия нормировки выполняется элементарным делением на длину, а для матрицы направляющих косинусов требуется более сложная процедура ортогонализации, что увеличивает вычислительную нагрузку.
\end{enumerate}


\subsection{Алгебра кватернионов}

Кватернион $q$ - гиперкомплексное число, которое может быть представлено в следующем виде \cite{hypercomplex_number, salamin1974quaternions}:
\begin{equation*}
	q = q_{\omega} + q_x i + q_y j + q_z k \text{,}
\end{equation*}
где $q_x, q_y, q_z$ - действительные числа, а $i, j, k$ - мнимые единицы, подчиняющиеся правилам умножения Гамильтона:
\begin{align*}
	i^2 = j^2 =& \hspace*{5pt} k^2 = ijk = -1 \\
	ij = k  \hspace*{15pt} jk &= i  \hspace*{15pt} ki = j 	\\
	ji = -k \hspace*{15pt} kj &= -i \hspace*{15pt} ik = -j
\end{align*}

Для удобства часто используются векторно-скалярная форма записи кватерниона: $q = (q_{\omega}, \bm{q_v})$, где $\bm{q_v} = (q_x, q_y, q_z)$.


\paragraph{Сложение кватернионов}
Сложение кватернионов выполняется покомпонентно \cite{hypercomplex_number}:
\begin{equation*}
	p + q = (p_{\omega} + q_{\omega}) + (p_x + q_x) i + (p_y + q_y) j + (p_z + q_z) k \text{.}
\end{equation*}


\paragraph{Умножение кватернионов}
Пусть даны два кватерниона $p = (p_w, p_x, p_y, p_z)$ и $q = (q_w, q_x, q_y, q_z)$. Их произведение $r = p \otimes q$ вычисляется по правилу \cite{hypercomplex_number}:

\begin{equation}
	\label{eq:q_mult}
	r = \begin{pmatrix}
		p_w q_w - p_x q_x - p_y q_y - p_z q_z \\
		p_w q_x + p_x q_w + p_y q_z - p_z q_y \\
		p_w q_y - p_x q_z + p_y q_w + p_z q_x \\
		p_w q_z + p_x q_y - p_y q_x + p_z q_w
	\end{pmatrix}^T \text{.}
\end{equation}

Умножение кватернионов ассоциативно, но не коммутативно \cite{hypercomplex_number}:
\begin{align*}
	(p \otimes q) \otimes r &= p \otimes (q \otimes r) 	\\
	p \otimes q &\neq q \otimes p
\end{align*}
Это соответствует некоммутативности композиции поворотов в трёхмерном пространстве.


\paragraph{Норма кватерниона}
Норма кватерниона вычисляется как:
\begin{equation*}
	\norm{q} = \sqrt{q_w^2 + q_x^2 + q_y^2 + q_z^2} \text{.}
\end{equation*}

Кватернион с единичной нормой $\norm{q} = 1$ называется единичным кватернионом. Его можно представить в следующем виде:
\begin{equation}
	\label{eq:geometry_interp_q}
	q = \left( \cos \frac{\theta}{2}, \bm{u} \sin \frac{\theta}{2} \right) \text{.}
\end{equation}
Такой кватернион описывает поворот в трёхмерном пространстве на угол $\theta$ вокруг оси, направление которой задаётся единичным вектором $\bm{u}$.

Процедура нормализации (перенормировки) выполняется тривиальным делением кватерниона на его норму:
\begin{equation}
	\label{eq:q_norm}
	q_{norm} = \frac{q}{\norm{q}} \text{.}
\end{equation}

Это одно из ключевых преимуществ кватернионов перед матрицами поворота, где требуется более сложная процедура ортогонализации.


\paragraph{Обратный и сопряжённый кватернионы}
Для кватерниона $q = (q_{\omega}, q_x, q_y, q_z)$ сопряжённый кватернион определяется как $q^* = (q_{\omega}, -q_x, -q_y, -q_z)$. Кватернион, обратный по умножению к $q$, вычисляется следующим образом:
\begin{equation*}
	q^{-1} = \frac{q^*}{\norm{q}^2} \text{.}
\end{equation*}


\paragraph{Поворот вектора}
Пусть задан вектор $\bm{v}$ в трёхмерном пространстве. Соответствующий ему <<чистый>> кватернион (кватернион с нулевой скалярной частью) имеет вид $q_v = (0, \bm{v})$. Тогда повёрнутый вектор $\bm{v'}$ получается с помощью двойного произведения кватернионов \cite{salamin1974quaternions, lee1991stereo}:
\begin{equation*}
	q_{v'} = (0, \bm{v'}) = q \otimes (0, \bm{v}) \otimes q^{-1} |_{\norm{q} = 1} = q \otimes (0, \bm{v}) \otimes q^{*} \text{.}
\end{equation*}
Результатом является кватернион с нулевой скалярной частью, векторная часть которого представляет собой искомый повёрнутый вектор.

В дальнейшем будем использовать следующую компактную запись вращения вектора с помощью кватерниона:
\begin{equation}
	\label{eq:vec_rotation}
	\bm{v'} = q \otimes (0, \bm{v}) \otimes q^{*} \text{.}
\end{equation}


\paragraph{Пример}
Пусть требуется повернуть вектор $\bm{v} = (1, 0, 0)$ на угол $\theta = 90 ^\circ$ вокруг оси $\bm{u} = (0, 0, 1)$ (ось $OZ$). Соответствующий кватернион вращения, согласно \eqref{eq:geometry_interp_q}:
\begin{equation*}
	q = \left( \cos(45 ^\circ), (0, 0, 1) \sin(45 ^\circ) \right) = \left( \frac{\sqrt{2}}{2}, 0, 0, \frac{\sqrt{2}}{2} \right) \text{.}
\end{equation*}

Выполняя вращение вектора $\bm{v}$ с помощью \eqref{eq:vec_rotation} и \eqref{eq:q_mult}, получим вектор $\bm{v'} = (0, 1, 0)$, что соответствует ожидаемому результату поворота.


\subsection{Применение кватернионов для задач навигации}

Широкий интерес механиков и специалистов в области навигации и управления движением к кватернионам возник начиная с 60-х годов прошлого века в связи с созданием автономных высокоточных компьютеризированных систем ориентации космических аппаратов, ракет, самолётов и морских кораблей. Разработка таких систем потребовала новых регулярных моделей, не содержащих особых точек типа деления на ноль, которые могли бы эффективно реализовываться в реальном времени на бортовых компьютерах. Кватернионные методы заняли здесь прочное место как теоретический аппарат, позволяющий наиболее эффективно решать задачи навигации.


\paragraph{Преобразование кватерниона в матрицу поворота}
Пусть задан единичный кватернион $q = (q_{\omega}, q_x, q_y, q_z)$, описывающий поворот от связанной системы координат к навигационной. Соответствующая матрица поворота $R$ вычисляется следующим образом \cite{salamin1974quaternions}:

\begin{equation}
	\label{eq:rotation_matrix_from_q}
	\bm{R} =
	\begin{bmatrix}
		1 - 2(q_x^2 + q_z^2) & 2(q_x q_y - q_w q_z) & 2(q_x q_z + q_w q_y) \\
		2(q_x q_y + q_w q_z) & 1 - 2(q_x^2 + q_z^2) & 2(q_y q_z - q_w q_x) \\
		2(q_x q_z - q_w q_y) & 2(q_y q_z + q_w q_x) & 1 - 2(q_x^2 + q_y^2)
	\end{bmatrix} \text{.}
\end{equation}

Эта формула непосредственно следует из правила умножения кватернионов и обеспечивает ортогональность матрицы при условии единичной нормы исходного кватерниона.


\paragraph{Преобразование матрицы поворота в кватернион}
Обратная задача — восстановление кватерниона по известной ортогональной матрице $R = \{r_{ij}\}$. Простейший подход основан на вычислении следа матрицы \cite{salamin1974quaternions}:
\begin{equation*}
	\operatorname{tr}(R) = r_{11} + r_{22} + r_{33} = 4q_w^2 - 1 \text{.}
\end{equation*}

Отсюда можно получить $q_w = \frac{1}{2} \sqrt{1 + \operatorname{tr}(R)}$. Остальные компоненты находятся из соотношений:
\begin{equation*}
	q_x = \frac{r_{32} - r_{23}}{4q_w}, \qquad
	q_y = \frac{r_{13} - r_{31}}{4q_w}, \qquad
	q_z = \frac{r_{21} - r_{12}}{4q_w} \text{.}
\end{equation*}

Данный метод работает корректно при $q_w \neq 0$. Однако при малых углах поворота $q_w$ может оказаться близким к нулю, что приводит к численной неустойчивости.


\paragraph{Кинематическое уравнение вращения в кватернионах}
Центральным соотношением, связывающим изменение ориентации тела с показаниями гироскопов, является кватернионное кинематическое уравнение. Если кватернион $q(t)$ задаёт ориентацию связанной системы координат относительно навигационной в момент времени $t$, а $\bm{\omega}(t) = \omega_x + \omega_y + \omega_z$ — вектор угловой скорости, измеряемой гироскопами, то скорость изменения кватерниона описывается уравнением \cite{lurie1961analiticheskaya}:
\begin{equation}
	\label{eq:q_kinematic_eq}
	\dot{q} = \frac{1}{2} q(t) \otimes (0, \bm{\omega}(t)) \text{.}
\end{equation}

Данное уравнение можно решать различными способами интегрирования, но если можно предположить, что угловая скорость $\bm{\omega}$ постоянна на всём промежутке $\Delta t$, то решение данного уравнения принимает следующий вид \cite{lurie1961analiticheskaya}:
\begin{equation}
	\label{eq:delta_q}
	q(t + \Delta t) = \left( \cos \frac{\norm{\bm{\omega}} \Delta t}{2}, \frac{\bm{\omega}}{\norm{\bm{\omega}}} \sin \frac{\norm{\bm{\omega}} \Delta t}{2} \right) \text{.}
\end{equation}

Важной особенностью численного интегрирования кватернионных уравнений является накопление ошибок округления, приводящих к отклонению нормы кватерниона от единицы. Поскольку только единичный кватернион соответствует повороту в трёхмерном пространстве, необходимо выполнять перенормировку согласно \eqref{eq:q_norm}.

% -----------------------------------------------------------------------------

\section{Начальная выставка (инициализация) по статическому буферу}
\label{sec:initial_alignment}

Перед началом движения путеизмерительная тележка в течение некоторого времени находится в неподвижном состоянии. В этот период записывается так называемый статический буфер — массив измерений акселерометра и гироскопа. Обработка этих данных позволяет решить две важнейшие задачи:
\begin{enumerate}
	\item определить начальную ориентацию связанной системы \(B\) относительно навигационной системы \(N\) (кватернион \(q_{body}^{nav}(0)\));
	\item вычислить смещения нуля акселерометра \(\bm{b}_{acc}\) и гироскопа \(\bm{b}_{gyro}\), которые затем будут вычитаться из измерений на протяжении всего проезда.
\end{enumerate}

В состоянии покоя акселерометр измеряет сумму проекции вектора гравитации \(\bm{g}\) на оси чувствительности, собственного смещения нуля $\bm{b}_{acc}$ и шумовой составляющей $\bm{\eta}_{acc}$:
\begin{equation}
	\bm{a}_{meas} = \bm{g}_{body} + \bm{b}_{acc} + \boldsymbol{\eta}_{acc} \text{.}
\end{equation}

Аналогично, гироскоп измеряет сумму проекцию угловой скорости вращения Земли $\omega_{earth}^{body}$, собственного смещения нуля $\bm{b}_{gyro}$ и шумовой составляющей $\bm{\eta}_{gyro}$:
\begin{equation}
	\bm{\omega}_{meas} = \bm{\omega}_{earth}^{body} + \bm{b}_{gyro} + \bm{\eta}_{gyro} \text{.}
\end{equation}


% -----------------------------------------------------------------------------

\subsection{Определения первоначальной ориентации тележки}

Для определения начальной ориентации тележки необходимо найти матрицу перехода $R_{B}^{N}$ из связанной системы координат $B$ в навигационную систему $N$. Эта матрица позволяет в дальнейшем вычислить соответствующий кватернион вращения. Построение $R_{B}^{N}$ основано на использовании трёх линейно независимых векторов, координаты которых известны в обеих системах. Выражая базисные векторы одной системы через координаты этих векторов в другой, можно непосредственно сформировать искомую матрицу перехода.

В качестве таких опорных векторов выступают:
\begin{itemize}
	\item вектор гравитации $\bm{g}$, измеряемый акселерометром в связанной системе и точно известный в навигационной системе $\bm{g}_{N}=(0,0,g)$, где $g = 9{,}8156 \text{ м}/\text{с}^2$ для широты Москвы;
	\item вектор угловой скорости вращения Земли $\bm{\omega}_{\text{earth}}$, измеряемый гироскопом в связанной системе и также известный в навигационной системе $\bm{\omega}_{\text{earth}} = (0, \omega_{e} \cos(L), \omega_{e} \sin(L))$, где $\omega_{e} = 7{,}2921 \times 10^{-5}$ рад/с, $L$ - географическая широта измерения;
	\item их векторное произведение $\bm{a} = \bm{\omega}_{\text{earth}} \times \bm{g}$, которое заведомо неколлинеарно с первыми двумя, что гарантирует линейную независимость тройки векторов.
\end{itemize}


В системе координат ENU данные векторы имеют следующие координаты:
\begin{align*}
	\bm{g} &= g * \bm{k} \\
	\bm{\omega} &= \omega_{y} * \bm{j} + \omega_{z} * \bm{k} \\
	\bm{a} &= a * \bm{i} \text{,\ \ \ \ где $a = \left| \left[ \bm{\omega}_{\text{earth}} \times \bm{g} \right] \right| $}
\end{align*}


А в системе координат тележки:
\begin{align*}
	\bm{g}' &= g_{x}' * \bm{i}' + g_{y}' * \bm{j}' + g_{z}' * \bm{k}' \\
	\bm{\omega}' &= \omega_{x}' * \bm{i}' + \omega_{y}' * \bm{j}' + \omega_{z}' * \bm{k}' \\
	\bm{a}' &= a_{x}' * \bm{i}' + a_{y}' * \bm{j}' + a_{z}' * \bm{k}' \\
\end{align*}


Получаем следующую систему уравнений:
\[
\begin{cases}
	\begin{aligned}
		&g * \bm{k} = g_{x}' * \bm{i}' + g_{y}' * \bm{j}' + g_{z}' * \bm{k}' \\
		&\omega_{y} * \bm{j} + \omega_{z} * \bm{k} = \omega_{x}' * \bm{i}' + \omega_{y}' * \bm{j}' + \omega_{z}' * \bm{k}' \\
		&a * \bm{i} = a_{x}' * \bm{i}' + a_{y}' * \bm{j}' + a_{z}' * \bm{k}' \\
	\end{aligned}
\end{cases}
\] \\


Из этой системы выразим $\bm{i}, \bm{j}, \bm{k}$:
\[
\begin{cases}
	\begin{aligned}
		\bm{i} &= \frac{a_{x}'}{a} * \bm{i}' + \frac{a_{y}'}{a} * \bm{j}' + \frac{a_{z}'}{a} * \bm{k}' \\
		\bm{j} &= \frac{1}{\omega_{y}}
		\left(
		\left[ \omega_{x}' - \omega_{z} \frac{g_{x}'}{g} \right] * \bm{i}' +
		\left[ \omega_{y}' - \omega_{z} \frac{g_{y}'}{g} \right] * \bm{j}' +
		\left[ \omega_{z}' - \omega_{z} \frac{g_{z}'}{g} \right] * \bm{k}'
		\right) \\
		\bm{k} &= \frac{g_{x}'}{g} * \bm{i}' + \frac{g_{y}'}{g} * \bm{j}' + \frac{g_{z}'}{g} * \bm{k}'
	\end{aligned}
\end{cases}
\] \\


Коэффициенты разложения базисных векторов $i, j, k$ по базису $i', j', k'$ образуют матрицу перехода из навигационной системы в связанную $R_N^B$, описывающую поворот из навигационной системы в связанную. По этой матрице вычисляем соответствующий единичный кватернион $q_N^B$. Тогда искомый кватернион ориентации тела относительно навигационной системы, т.е. преобразование из связанной системы в навигационную, находится как сопряжённый: $q_{body}^{nav} = \left( q_N^B \right)^{-1} = \left( q_N^B \right)^{*}$. Такой подход позволяет избежать дополнительного обращения матрицы и обеспечивает численную устойчивость за счёт внутренней ортогонализации при преобразовании матрицы в кватернион.


\subsection{Определения смещений чувствительных осей датчиков}

Для определения смещений акселерометра и гироскопа из собранных статических значений $\bm{\bar{a}}$ и $\bm{\bar{\omega}}$ необходимо вычесть проекции $\bm{g}_{N}$ и $\bm{\omega}_{\text{earth}}$ на оси датчиков в начальный момент:

\begin{equation*}
	\bm{b}_{acc} = \bm{\bar{a}} - \bm{g}_{earth}^{\hspace{2pt} body}				\hspace{40pt}
	\bm{b}_{gyro} = \bm{\bar{\omega}} - \bm{\omega}_{earth}^{\hspace{2pt} body}
\end{equation*}

В соответствии с \eqref{eq:vec_rotation}:
\begin{align*}
	\bm{g}_{earth}^{\hspace{2pt} body} &= \left( q_{body}^{nav} \right)^{-1} \otimes
	\left( 0, \bm{g}_{earth}^{\hspace{2pt} nav} \right) \otimes q_{body}^{nav}	\\
	\bm{\omega}_{earth}^{body} &= \left( q_{body}^{nav} \right)^{-1} \otimes
	\left( 0, \bm{\omega}_{earth}^{nav} \right) \otimes q_{body}^{nav}
\end{align*}

Итоговые выражения для смещений чувствительных осей акселерометра и гироскопа:
\begin{align}
	\bm{b}_{acc} &= \bm{\bar{a}} - \left( q_{body}^{nav} \right)^{-1} \otimes
	\left( 0, \bm{g}_{earth}^{\hspace{2pt} nav} \right) \otimes q_{body}^{nav}		\label{eq:bias_acc} \\
	\bm{b}_{gyro} &= \bm{\bar{\omega}} - \left( q_{body}^{nav} \right)^{-1} \otimes
	\left( 0, \bm{\omega}_{earth}^{nav} \right) \otimes q_{body}^{nav}				\label{eq:bias_gyro}
\end{align}


% -----------------------------------------------------------------------------

\section{Алгоритм построения траектории движения}
\label{sec:trajectory}

В данном разделе будет описан алгоритм построения траектории движения по данным с гироскопа и акселерометра.

\subsection{Получение ориентации тела в каждый момент времени}
Первая задача для построения траектории - получить точную оценку ориентации тела (путеизмерительной тележки) относительно навигационной системы координат (ENU) в каждый момент времени. Ориентация представляется кватернионом единичной нормы $q_{body}^{nav}(t_i)$.

Для начального момента времени $q_{body}^{nav}(t_0)$ совпадает с кватернионом, полученным при определении первоначальной ориентации тележки.

В движении измерения гироскопа представляют собой следующее выражение (без учёта шума):

\begin{equation*}
	\bm{\omega_{meas}}(t_i) = \bm{\omega}_{body}(t_i) + \bm{b}_{gyro} + \bm{\omega}_{earth}^{body}(t_i)
\end{equation*}

где:
\begin{itemize}
	\item $\bm{\omega}_{body}(t_i)$ - истинная угловая скорость тела в момент времени $t_i$;
	\item $\bm{b}_{gyro}$ - смещение гироскопа;
	\item $\bm{\omega}_{earth}^{body}$ - проекция угловой скорости вращения Земли на оси гироскопа, вычисляемая как $\bm{\omega}_{earth}^{body}(t_i) = \left( q_{body}^{nav}(t_i) \right)^{-1} \otimes \left( 0, \bm{\omega}_{earth}^{nav} \right) \otimes q_{body}^{nav}(t_i)$.
\end{itemize}


Для перехода от $t_i$ к $t_{i+1}$ с шагом $\Delta t = t_{i+1} - t_i$ нужно вычислить приращение ориентации. Считая угловую скорость постоянной на интервале $\Delta t$, кватернион приращения согласно \eqref{eq:delta_q} будет равен:

\begin{equation}
	\Delta q = \left(
	\cos \frac{\theta}{2}, \bm{u} \sin \frac{\theta}{2}
	\right) \text{,}
\end{equation}

где $\theta = \norm{\bm{\omega}_{body}(t_i)} \Delta t$ - угол поворота, а $\bm{u} = \bm{\omega}_{body}(t_i) / \norm{\bm{\omega}_{body}(t_i)}$ - ось вращения.

Тогда кватернион для описания ориентации тележки при $t = t_{i+1}$, будет равен:

\begin{equation}
	q_{body}^{nav}(t_{i+1}) = \Delta q \otimes q_{body}^{nav}(t_i) \text{, }
\end{equation}

что соответствует последовательным поворотам сначала из ориентации при $t = t_{i+1}$ в i-ую ориентацию, а затем в навигационную систему координат.

После умножения норма кватерниона может отклониться от единицы из-за численных ошибок, поэтому необходимо провести его нормировку:

\begin{equation*}
	q_{body}^{nav}(t_{i+1}) = \frac{q_{body}^{nav}(t_{i+1})}{\norm{q_{body}^{nav}(t_{i+1})}} \text{.}
\end{equation*}

Таким образом получен массив кватернионов для описания ориентации в каждый момент времени измерения.


\subsection{Обработка данных акселерометра}
Вторым шагом для определения траектории движения путеизмерительной тележки является необходимость получения вектора линейного ускорения в навигационной системе координат, свободный от влияния гравитации и инструментальных погрешностей датчика. Исходные измерения акселерометра представляют собой проекции кажущегося ускорения на оси чувствительности, то есть сумму ускорения движения тележки, проекции вектора гравитации и постоянного смещения нуля датчика:
\begin{equation}
	\bm{a}_{meas}(t_i) = \bm{a}_{body}(t_i) + \bm{b}_{acc} + \bm{g}_{earth}^{body}(t_i)
\end{equation}

Задача обработки заключается в последовательном устранении этих составляющих.

\paragraph{Коррекция смещения акселерометра}
Первым этапом из измеренного сигнала вычитается постоянное смещение, определённое в процессе начальной выставки. Поскольку смещение присуще датчику в его собственной системе координат, эта операция выполняется до преобразования системы отсчёта:
\begin{equation*}
	\bm{a}_{corr}(t_i) = \bm{a}_{meas}(t_i) - \bm{b}_{acc}
\end{equation*}

\paragraph{Поворот ускорения в навигационную систему координат и вычитание гравитации}
Полученный вектор $a_{corr}$ задан в системе координат тележки, так что для его перевода в навигационную систему ENU используется кватернион ориентации $q_{body}^{nav}(t_i)$ в соответствии с \eqref{eq:vec_rotation}:
\begin{equation*}
	\bm{a}_{nav}(t_i) = q_{body}^{nav}(t_i) \otimes (0, \bm{a}_{corr}(t_i)) \otimes (q_{body}^{nav}(t_i))^{-1} \text{.}
\end{equation*}

В результате получается вектор $\bm{a}_{nav}(t_i)$ - кажущееся ускорение, выраженное в навигационной системе координат.

Последним этапом обработки данных акселерометра является вычитание постоянного вектора гравитации из вектора $\bm{a}_{nav}(t_i)$:
\begin{equation*}
	\bm{a}_{corr}^{nav}(t_i) = \bm{a}_{nav}(t_i) - \bm{g_{nav}} \text{.}
\end{equation*}

В конечном итоге, получен массив ускорения путеизмерительной тележки в навигационной системе координат ENU. \\


\subsection{Расчёт скорости и построение траектории движения}
После получения массива ускорений движения $\bm{a}_{corr}^{nav}(t_i)$ в навигационной системе координат следующим этапом является численное интегрирование для вычисления скорости и положения тележки. Ввиду неизбежного накопления инструментальных погрешностей интегрирование дополняется алгоритмами коррекции, использующими априорную информацию о движении.

\paragraph{Интегрирование ускорений для получения скорости}
Априорной информацией для коррекции скорости является факт полной остановки в конце проезда, а также неподвижность тележки во время выставки датчиков перед началом проезда, т.е. $\bm{v}(t_0) = \bm{v}_{real}(t_N) = 0$.

Для интегрирования ускорений используется метод трапеций, обеспечивающий второй порядок точности вычислений. Для каждого момента времени $t_i = t_0, ..., t_N$ с шагом в $\Delta t_i = t_i - t_{i-1}$ вычисление скорости выполняется по следующей формуле:
\begin{equation}
	\bm{v}(t_i) = \bm{v}(t_{i-1}) + \frac{\bm{a}_{corr}^{nav}(t_{i-1}) + \bm{a}_{corr}^{nav}(t_i)}{2} \Delta t_i \text{.}
\end{equation}

В результате получается массив скорости $\bm{v}(t_i)$, который уже содержит накопленную ошибку интегрирования. Предполагая, что основная составляющая ошибки ускорения постоянна во времени, ошибка скорости накапливается линейно. Поэтому поправка вводится линейно убывающей от нуля в начале до конечного значения, равное $\Delta \bm{v} = \bm{v}(t_N) - \bm{v}_{real}(t_N) = \bm{v}(t_N)$.

Таким образом, скорректированное значение скорости вычисляется следующим соотношением:
\begin{equation}
	\bm{v}_{corr}(t_i) = \bm{v}(t_i) - \Delta \bm{v} * \frac{t_i - t_0}{t_N - t_0} \text{.}
\end{equation}


\paragraph{Расчёт перемещения по скорректированной скорости}
После получения скорректированной скорости $\bm{v}(t_i)$ вычисляется перемещение $\bm{p}(t_i)$ тележки относительно первоначального положения, при котором $\bm{p}(t_0) = (0, 0, 0)$. Интегрирование скорости также проводится с помощью метода трапеций:
\begin{equation}
	\bm{p}(t_i) = \bm{p}(t_{i-1}) + \frac{\bm{v}_{corr}(t_{i-1}) + \bm{v}_{corr}(t_i)}{2} \Delta t_i \text{.}
\end{equation}

Начальное и конечное положение тележки априорно известны, например, данные широты и долготы могут быть получены по спутниковому приёмнику, а перепад высот может быть измерен высотомером, так что перемещение по каждой оси системы координат ENU может быть заранее вычислено. Тогда ошибка перемещений в конце может быть вычислена по формуле:
\begin{equation*}
	\Delta \bm{p} = \bm{p}(t_N) - \bm{p}_{real} \text{.}
\end{equation*}

Аналогично скорости, ошибку перемещения корректируем линейно:
\begin{equation}
	\bm{p}_{corr}(t_i) = \bm{p}(t_i) - \Delta \bm{p} * \frac{t_i - t_0}{t_N - t_0} \text{.}
\end{equation}

Таким образом построена траектория движения путеизмерительной тележки на основе зашумлённых данных инерциальных датчиков, в соответствии с изначально поставленной задачей.